% Generated draft: requires mathematical and source review.
\documentclass[11pt]{article}
\usepackage[margin=25mm]{geometry}
\usepackage{fontspec}
\setmainfont{Latin Modern Roman}
\usepackage{amsmath,amssymb,mathtools,unicode-math}
\setmathfont{Latin Modern Math}
\usepackage{enumitem,needspace,xurl,hyperref}
\setlength{\parindent}{0pt}
\setlength{\parskip}{5pt}
\setlength{\emergencystretch}{3em}
\newcommand{\N}{\mathbb{N}}
\newcommand{\Z}{\mathbb{Z}}
\newcommand{\Q}{\mathbb{Q}}
\newcommand{\R}{\mathbb{R}}
\newcommand{\C}{\mathbb{C}}
\newcommand{\F}{\mathbb{F}}
\newcommand{\A}{\mathbb{A}}
\newcommand{\PP}{\mathbb{P}}
\newcommand{\E}{\mathbb{E}}
\newcommand{\eps}{\varepsilon}
\newcommand{\dd}{\,\mathrm{d}}
\begin{document}
\section*{UnsolvedMath ID: 30001174}
\textbf{Source problem code:} catalogue entry
\subsection{Definitions and mathematical statement}
Category theory studies objects and morphisms, with associative composition of composable morphisms and an identity morphism for each object. The phrase ``natural use'' below is an informal methodological judgment; the supplied record does not identify it with the technical notion of a natural transformation.

\begin{quote}
Should category theory be applied everywhere in mathematics, or are there mathematical questions or methods for which its use is not natural?
\end{quote}

This is the supplied methodological question. It does not specify a domain of mathematical questions or methods, what constitutes an application of category theory, or a criterion for such an application to be natural. Moreover, ``should be applied'' expresses a methodological recommendation rather than a defined mathematical predicate. Consequently, the supplied evidence does not determine a precise quantified mathematical proposition. Turning this question into a theorem or conjecture would require additional choices that are not supplied here.
\subsection{Short English statement}
The question asks whether categorical methods are appropriate throughout mathematics or have natural limits. The supplied record leaves appropriateness and the scope of application undefined, so this remains a methodological question requiring clarification before mathematical formalization.
\subsection{Sources}
\begin{itemize}
\item Supplied local UnsolvedMath registry record, permanent ID \(30001174\), title ``Natural Limits of Categorical Methods,'' source problem code \texttt{catalogue entry}, in the collection ``Oberwolfach Reports Open Problems.'' The record gives the proposed year as 2009 and names no proposer.
\item Supplied report citation: \emph{Mini-Workshop: Category Theory and Related Fields: History and Prospects} (2009), Oberwolfach Reports, pp.~463--492, DOI 10.4171/owr/2009/08. \url{https://doi.org/10.4171/owr/2009/08}. The background describes Oberwolfach Reports as documenting mathematical workshops and their open-problem discussions. The supplied literature-source label is ``OWR Category Theory and Related Fields,'' with the same URL.
\item The registry labels the statement ``exact'' and reports that its mathematical content was checked against the official report, with notation, context, and question form normalized without changing the source-level claim. This is supplied verification metadata; the report has not been independently consulted for this draft.
\item Supplied literature assessment, checked on \texttt{2026-08-21}: status ``open.'' The assessment describes the question as a broad methodological question rather than a settled mathematical proposition, and as remaining open-ended. This is supplied dated curation, not an independently verified literature assessment. The supplied background requests independent verification of time-sensitive status and cited claims.
\end{itemize}
\end{document}
